% ============================================================
% §2  Pole Φ — źródła, superpozycja, interferencja
% ============================================================
\section{Pole generuj\k{a}ce przestrze\'n}\label{sec:pole}

\subsection{Definicja pola $\Phi$}

\begin{definition}[Pole przestrzenno\'sci $\Phi$]\label{def:Phi}
Pole $\Phi$ jest rzeczywistym polem skalarnym, kt\'orego warto\'s\'c
w~danym punkcie mierzy \textbf{g\k{e}sto\'s\'c wygenerowanej przestrzeni}
w~tym miejscu. Nie jest ono polem na gotowej rozmaito\'sci ---
jest polem, kt\'ore t\k{e} rozmaito\'s\'c \emph{konstytuuje}:
\begin{itemize}[nosep]
  \item $\Phi = 0$: brak przestrzeni (nico\'s\'c $\Nzero$),
  \item $\Phi > 0$: przestrze\'n istnieje, jej ,,g\k{e}sto\'s\'c''
        wynosi~$\Phi$,
  \item $\Phi \to \infty$: ekstremalnie g\k{e}sta przestrze\'n
        (wn\k{e}trze obiekt\'ow masywnych).
\end{itemize}
\end{definition}

Pole $\Phi$ nie jest \textbf{dodatkowym polem} w~przestrzeni.
Jest polem, \textbf{z~kt\'orego} przestrze\'n si\k{e} sk\l{}ada.
Metryka, odleg\l{}o\'s\'c, przyczynowo\'s\'c --- to w\l{}asno\'sci
konfiguracji $\Phi$, nie oddzielne byty.

\begin{remark}[Interpretacja fizyczna]
Mo\.zna my\'sle\'c o~$\Phi$ jak o~,,g\k{e}sto\'sci eteru'' ---
ale nie eteru jako o\'srodka materialnego, lecz jako samej
\emph{zdolno\'sci do istnienia relacji przestrzennych}.
Tam, gdzie $\Phi = 0$, nie ma niczego --- nawet pustki.
Tam, gdzie $\Phi > 0$, obiekty mog\k{a} wzgl\k{e}dem siebie
istnie\'c, oddzia\l{}ywa\'c, porusza\'c si\k{e}.
\end{remark}

% ----------------------------------------------------------
\subsection{Materia jako \'zr\'od\l{}o pola $\Phi$}\label{ssec:zrodla}

Ka\.zda cz\k{a}stka fundamentalna generuje wk\l{}ad do pola~$\Phi$.
Nie jest to sprzężenie ,,zewn\k{e}trzne'' (materia w~przestrzeni
sprz\k{e}\.zona z~polem), lecz \textbf{akt konstytutywny}: cz\k{a}stka
\emph{jest} lokalnym wzbudzeniem, kt\'ore jednocze\'snie wytwarza
pole przestrzenno\'sci wok\'o\l{} siebie.

\begin{axiom}[Materia generuje przestrze\'n]\label{ax:zrodlo}
Ka\.zde \'zr\'od\l{}o masy--energii $\rho_a$ wytwarza wk\l{}ad do pola $\Phi$:
\begin{equation}\label{eq:Phi-source}
    \mathcal{D}[\Phi_a] = q_a \cdot \rho_a,
\end{equation}
gdzie $q_a > 0$ jest \textbf{sta\l{}\k{a} generacji przestrzennej}
cz\k{a}stki typu~$a$ (analogiczna do \l{}adunku, ale dla przestrzeni),
a~$\mathcal{D}$ jest operatorem opisuj\k{a}cym dynamik\k{e} pola $\Phi$.
\end{axiom}

\begin{remark}[Sta\l{}a generacji $q_a$]
Sta\l{}a $q_a$ okre\'sla, jak silnie dany typ cz\k{a}stki generuje
przestrze\'n. Mo\.ze zale\.ze\'c od masy, spinu, lub innych w\l{}asno\'sci
cz\k{a}stki. Jest to nowy parametr fizyczny, w\l{}a\'sciwy TGP,
nieobecny w~standardowej fizyce. W~granicy jednorodnej
($q_a = q$ dla wszystkich typ\'ow) teoria upraszcza si\k{e}
do jednoparametrowej.
\end{remark}

\begin{remark}[Natura operatora $\mathcal{D}$]\label{rem:D-natura}
Operator $\mathcal{D}$ w~r\'ownaniu~\eqref{eq:Phi-source} nie jest
standardowym laplasjanem na gotowej rozmaito\'sci --- to
narusza\l{}oby aksjomat~\ref{ax:przestrzen}. Musi by\'c
\textbf{wewn\k{e}trzny wzgl\k{e}dem $\Phi$}: zdefiniowany
przez sam\k{a} konfiguracj\k{e} pola konstytuuj\k{a}cego
przestrze\'n. Jego jawna posta\'c jest wyprowadzona
w~sekcji~\ref{sec:formalizm} (definicja~\ref{def:D}).

Dla referencji: cz\k{e}\'s\'c liniowa operatora wynosi
\begin{equation}\label{eq:D-preview}
  \mathcal{D}[\Phi]\big|_{\mathrm{lin}}
  = \frac{1}{\PhiZero}\,\nabla^2\Phi,
\end{equation}
a~pe\l{}ne r\'ownanie pola zawiera dodatkowo nieliniowy cz\l{}on
samointerferencji $\mathcal{N}[\Phi]$ (definicja~\ref{def:N},
sekcja~\ref{sec:formalizm}):
\begin{equation}\label{eq:N-preview}
  \mathcal{N}[\Phi]
  = \frac{\alpha}{\PhiZero}\,\frac{(\nabla\Phi)^2}{\Phi}
  + \beta\,\frac{\Phi^2}{\PhiZero^2}
  - \gamma\,\frac{\Phi^3}{\PhiZero^3},
\end{equation}
gdzie $\alpha = 2$ (z~wariacji dzia\l{}ania),
$\beta, \gamma > 0$, $[\beta] = [\gamma] = \mathrm{L}^{-2}$.
\end{remark}

% ----------------------------------------------------------
\subsection{Pojedyncze \'zr\'od\l{}o: sferyczna
  przestrze\'n generowana}\label{ssec:pojedyncze}

Dla izolowanej, statycznej cz\k{a}stki punktowej o~masie $M$
w~punkcie $\mathbf{x}_0$ pole generowane ma symetri\k{e} sferyczn\k{a}:
\begin{equation}\label{eq:Phi-single}
    \Phi_{\mathrm{poj}}(r)
    = \frac{q\,M}{4\pi}\,\phi(r),
    \qquad
    r = |\mathbf{x} - \mathbf{x}_0|,
\end{equation}
gdzie $\phi(r)$ jest \textbf{profilem generacji} --- funkcj\k{a}
zale\.zn\k{a} od odleg\l{}o\'sci od \'zr\'od\l{}a.

Profil $\phi(r)$ \textbf{nie jest} z~g\'ory za\l{}o\.zony jako
$e^{-mr}/r$ (Yukawa) ani $1/r$ (Coulomb). Jego posta\'c wynika
z~pe\l{}nej, nieliniowej dynamiki pola $\Phi$ i~musi by\'c
\textbf{wyprowadzona} z~r\'owna\'n teorii. Na obecnym etapie
mo\.zemy jedynie za\l{}o\.zy\'c og\'olne w\l{}asno\'sci:
\begin{enumerate}[nosep]
  \item $\phi(r) > 0$ wsz\k{e}dzie (przestrze\'n jest generowana,
        nie ,,ujemna''),
  \item $\phi(r) \to \infty$ dla $r \to 0$ (w~centrum
        \'zr\'od\l{}a g\k{e}sto\'s\'c przestrzeni dywerguje),
  \item $\phi(r) \to 0$ dla $r \to \infty$ (daleko od \'zr\'od\l{}a
        przestrze\'n zanika),
  \item $\phi(r)$ jest niemonotoniczny --- posiada punkt przegi\k{e}cia
        lub lokalne ekstremum, co prowadzi do trzech re\.zim\'ow
        (sekcja~\ref{sec:rezimy}).
\end{enumerate}

% ----------------------------------------------------------
\subsection{Wiele \'zr\'ode\l{}: nak\l{}adanie
  i~interferencja}\label{ssec:interferencja}

W~obecno\'sci wielu cz\k{a}stek ca\l{}kowite pole $\Phi$ nie jest
prost\k{a} sum\k{a} wk\l{}ad\'ow poszczeg\'olnych \'zr\'ode\l{}.
Jest wynikiem \textbf{nieliniowej interferencji}.

\subsubsection{Przybli\.zenie liniowe (s\l{}abe pole)}

W~granicy s\l{}abego pola ($\Phi \ll \Phi_{\mathrm{char}}$, gdzie
$\Phi_{\mathrm{char}}$ jest skal\k{a} nieliniowo\'sci) wk\l{}ady
superponu\k{a} si\k{e} liniowo:
\begin{equation}\label{eq:Phi-linear}
    \Phi(\mathbf{x})
    \simeq \sum_{i} \Phi_i(\mathbf{x})
    = \sum_{i} \frac{q_i\,M_i}{4\pi}\,\phi(|\mathbf{x}-\mathbf{x}_i|).
\end{equation}
To przybli\.zenie odpowiada granicy newtonowskiej i~jest
wystarczaj\k{a}ce na du\.zych skalach.

\subsubsection{Pe\l{}ny re\.zim nieliniowy}

W~obszarach, gdzie wk\l{}ady od wielu \'zr\'ode\l{} nak\l{}adaj\k{a} si\k{e}
silnie (g\k{e}ste skupiska materii, wn\k{e}trza gwiazd, j\k{a}dra atom\'ow),
r\'ownanie pola staje si\k{e} \textbf{nieliniowe}:
\begin{equation}\label{eq:Phi-nonlinear}
    \mathcal{D}[\Phi]
    = \sum_{i} q_i \rho_i
    + \mathcal{N}[\Phi],
\end{equation}
gdzie $\mathcal{N}[\Phi]$ opisuje \textbf{samointerferencj\k{e}}
pola --- nieliniowe sprz\k{e}\.zenie wygenerowanej przestrzeni
samej z~sob\k{a}. Jawna posta\'c $\mathcal{N}[\Phi]$ podana jest
w~r\'ownaniu~\eqref{eq:N-preview} powy\.zej
oraz definicji~\ref{def:N} (sekcja~\ref{sec:formalizm});
jawna posta\'c pe\l{}nego r\'ownania pola ---
w~twierdzeniu~\ref{thm:field-eq}.

To w\l{}a\'snie cz\l{}on $\mathcal{N}[\Phi]$ jest odpowiedzialny za:
\begin{itemize}[nosep]
  \item odpychanie na \'srednich skalach (wysoki potencja\l{}
        przestrzenny mi\k{e}dzy bliskimi \'zr\'od\l{}ami),
  \item studni\k{e} na ekstremalnie ma\l{}ych skalach (sklejenie),
  \item wzmocnienia i~wyga\'szenia pola w~zale\.zno\'sci od
        wzgl\k{e}dnej konfiguracji \'zr\'ode\l{}.
\end{itemize}

\begin{remark}[Interferencja, nie superpozycja]
S\l{}owo ,,interferencja'' jest celowe. W~fizyce fal interferencja
oznacza, \.ze wynik na\l{}o\.zenia zale\.zy od \emph{fazy} --- mog\k{a}
wyst\k{a}pi\'c wzmocnienia i~wyga\'szenia. W~TGP analogiczn\k{a} rol\k{e}
pe\l{}ni \textbf{gradient pola} $\nabla\Phi$: dwa bliskie \'zr\'od\l{}a
generuj\k{a} mi\k{e}dzy sob\k{a} obszar o~\emph{wy\.zszym} $\Phi$ ni\.z ka\.zde
z~osobna (konstruktywne na\l{}o\.zenie); ale si\l{}a wypadkowa na obiekty
zale\.zy od \emph{gradientu} $\Phi$, kt\'ory mo\.ze by\'c skierowany
zar\'owno ku \'zr\'od\l{}u, jak i~od niego --- w~zale\.zno\'sci od skali.
To jest mechanizm trzech re\.zim\'ow (sekcja~\ref{sec:rezimy}).
\end{remark}

% ----------------------------------------------------------
\subsection{Pole $\Phi$ a~metryka}\label{ssec:metryka}

Metryka nie jest obiektem fundamentalnym. Jest \textbf{opisem
efektywnym} konfiguracji pola~$\Phi$. W~regionach, gdzie
$\Phi > 0$ i~zmienia si\k{e} g\l{}adko, mo\.zna zdefiniowa\'c
efektywn\k{a} metryk\k{e}:
\begin{equation}\label{eq:metric-from-Phi}
    g_{\mu\nu}^{\mathrm{eff}} = g_{\mu\nu}[\Phi, \partial\Phi, \ldots\,].
\end{equation}
Dok\l{}adna posta\'c tej relacji jest \textbf{otwartym problemem} TGP.
Na obecnym etapie mo\.zna jedynie stwierdzi\'c:
\begin{itemize}[nosep]
  \item metryka jest okre\'slona tylko tam, gdzie $\Phi > 0$;
  \item jednorodne $\Phi = \mathrm{const} > 0$ daje p\l{}ask\k{a}
        metryk\k{e} (Minkowski);
  \item gradienty $\Phi$ generuj\k{a} krzywizn\k{e} --- to jest
        grawitacja w~TGP;
  \item du\.ze $\Phi$ z~du\.zymi gradientami odpowiada silnemu
        polu grawitacyjnemu.
\end{itemize}

\begin{hypothesis}[Metryka minimalna]\label{hyp:metric}
W~najprostszym uj\k{e}ciu metryka efektywna ma posta\'c:
\begin{equation}\label{eq:metric-minimal}
    ds^2 = -\frac{c(\Phi)^2}{c_0^2}\,dt^2
    + \frac{\Phi}{\PhiZero}\,\delta_{ij}\,dx^i dx^j,
\end{equation}
gdzie $c(\Phi)$ jest lokaln\k{a} pr\k{e}dko\'sci\k{a} \'swiat\l{}a
(sekcja~\ref{sec:stale}), a~$\PhiZero$ warto\'sci\k{a} referencyjn\k{a}.
Cz\l{}on przestrzenny jest modulowany przez $\Phi/\PhiZero$ ---
wi\k{e}cej wygenerowanej przestrzeni oznacza wi\k{e}ksz\k{a}
,,obj\k{e}to\'s\'c'' dost\k{e}pn\k{a} w~danym regionie.
\end{hypothesis}

Ta forma metryki jest \textbf{hipotez\k{a} robocz\k{a}}, nie
aksjomatem. Jej sp\'ojno\'s\'c z~obserwacjami musi by\'c
zweryfikowana. Poprawiona posta\'c (eksponencjalna), daj\k{a}ca
poprawne parametry post-newtonowskie, jest wyprowadzona
w~sekcji~\ref{ssec:metryka-popr}.

\begin{remark}[Metryka a~$\Ffield$]\label{rem:metryka-F}
Metryka efektywna jest okre\'slona tylko w~$\mathcal{S}_1$
(sekcja~\ref{ssec:sektory}), czyli tam, gdzie $\Ffield > 0$
i~$\Phi > 0$. W~$\mathcal{S}_0$ ($\Ffield = 0$, $\Phi = 0$)
poj\k{e}cie metryki traci sens --- nie ma ,,odleg\l{}o\'sci''
w~nico\'sci. Metryka jest zatem \textbf{bytem emergentnym}
wzgl\k{e}dem pola $\Phi$ --- polem, kt\'ore konstytuuje
mo\.zliwo\'s\'c mierzenia odleg\l{}o\'sci, nie polem ,,na''
gotowej przestrzeni.
\end{remark}
